\emph{(1).} This clause is a corollary of Theorem \claimref{synthesis.rung-field-d-term-action-and-tree-masses-on-the-lift-locus}(4), by a two-line elimination. A term linear in $m$ added to $W$ would itself be a level (it shifts $\zeta$), so functions of $\mathsf{N}$ are the general level-free case. By Theorem \claimref{synthesis.rung-field-d-term-action-and-tree-masses-on-the-lift-locus}(4), a family point is critical for $V = V_{D} + W(\mathsf{N})$ exactly when $\tfrac{g^{2}}{4}m^{*} + \tfrac{g'^{2}}{4}(m^{*} - \xi) + W'(\mathsf{N}^{*})\tfrac{m^{*}}{\mathsf{N}^{*}} = 0$, and there both charged modes have $m_{\pm}^{2} = 2\bigl(\tfrac{g^{2}}{4}\mathsf{N}^{*} + W'(\mathsf{N}^{*})\bigr)$. For $m^{*} \neq 0$ the first equation gives $W'(\mathsf{N}^{*}) = -\tfrac{g^{2} + g'^{2}}{4}\mathsf{N}^{*} + \tfrac{g'^{2}}{4}\tfrac{\xi}{m^{*}}\mathsf{N}^{*}$, and substituting,
\[
m_{\pm}^{2} = \frac{g'^{2}\mathsf{N}^{*}}{2}\Bigl(\frac{\xi}{m^{*}} - 1\Bigr) = (M_{Z}^{2} - M_{W}^{2})\Bigl(\frac{\xi}{m^{*}} - 1\Bigr),
\]
since $M_{Z}^{2} - M_{W}^{2} = \tfrac12 g'^{2}\mathsf{N}^{*} > 0$. Positivity is $\xi/m^{*} > 1$; on the family $m^{*} = -H^{2}\sinh 2\eta^{*}$ (Theorem \claimref{synthesis.rung-field-on-the-lift-locus}). For $\xi = 0$ the bracket is $-1$. The argument uses only $W'(\mathsf{N}^{*})$, so it holds for every $W$, whatever constants it carries. For $W = 0$ the vacuum equation gives $\xi/m^{*} = (g^{2} + g'^{2})/g'^{2}$ and the formula returns the stored $m_{\pm}^{2} = M_{W}^{2}$. The check does not rely on the stored mass formulas: for $W = \tfrac32(\mathsf{N} - 7)^{2}$, for the cubic $W = \tfrac14(\mathsf{N} - 2)^{3} + 5\mathsf{N}$, and for a quadratic $W$ chosen so that $\xi = 0$, in the frames $\lambda^{*} = \tfrac12$ and $5$ at $H^{2} = 2$, $g = \tfrac23$, $g' = \tfrac12$, it solves the level from criticality, computes the Hessian of $V$ on $L$ in the charged plane by exact polynomial interpolation (seventeen nodes, enough for the cubic), divides by the kinetic coefficient $2\mathsf{N}(d)$, and verifies the formula and the negative value at $\xi = 0$.

\emph{(2).} Write a real-linear map of $\mathcal{R}_{\C} = \C^{4}$ as a complex $4 \times 4$ matrix $A$ in the coordinates $(\alpha, \beta)$ (a real-linear infinitesimal isometry of $Q$ is complex-linear: write it as $A_{1} + A_{2}$ with $A_{1}$ linear and $A_{2}$ antilinear; the condition $Q(c, Ac) = 0$ for all $c$ splits into a holomorphic and a mixed part, so $Q(c, A_{2}d) + Q(d, A_{2}c) = 0$ for all $c, d$, and replacing $d$ by $I_{4}d$ gives $Q(c, A_{2}d) = Q(d, A_{2}c)$, hence $Q(c, A_{2}d) = 0$ and $A_{2} = 0$ since $Q$ is nondegenerate; the computation is done over the 32 real coordinates of a complex matrix $A$). The conditions $A^{\dagger} + A = 0$ (preserving $G$, whose matrix is a multiple of the identity), $A^{\mathsf T}S_{Q} + S_{Q}A = 0$ with $S_{Q}$ the matrix of $Q$, and $[A, T_{j}] = 0$ for $j = 1, 2, 3$ are real-linear in the 32 real coordinates; their solution space, computed exactly, has dimension $3$, and $R_{1}$, $R_{2}$, $Y$ satisfy them and are independent. The brackets are verified on a basis. For the right generators, $R_{1}^{2} = -\tfrac14$, so $(1 - tR_{1})^{-1} = (1 + tR_{1})/(1 + s)$ and the Cayley transform $(1 + tR_{1})(1 - tR_{1})^{-1}$ is $C_{t}$; it preserves $G$ and $Q$ because $tR_{1}$ is in their isometry algebra and $1 - tR_{1}$ is invertible. On $c_{\lambda}$, $2R_{1}c_{\lambda} = ((0, \lambda), (0, -H^{2}/\lambda))$, so $C_{t}c_{\lambda}$ has $\alpha = (aH^{2}/\lambda, b\lambda)$, $\beta = (-a\lambda, -bH^{2}/\lambda)$ with $a = (1 - s)/(1 + s)$, $b = t/(1 + s)$, $a^{2} + b^{2} = 1$, and $m(C_{t}c_{\lambda}) = (a^{2} - b^{2})m(c_{\lambda})$, the stated factor. The moment-map identities are checked at five Gaussian-rational points (three generic, one off the family on $L$, two on the family); the sum for $\su(2)_{R}$ equals the stored sum for the weak $\su(2)$. For invariant functions: by Theorem \claimref{synthesis.rung-field-d-term-action-and-tree-masses-on-the-lift-locus}(1) a $U(2)$-invariant function on $L$ is $f(\mathsf{N}, m)$ on the orbit space $\mathsf{N}^{2} - m^{2} \geq H^{4}$; at fixed $\mathsf{N} \geq H^{2}$ the fibre is $|m| \leq \sqrt{\mathsf{N}^{2} - H^{4}}$, which is $|m(c_{\lambda})|$ for the family point with that norm; the factor $a^{2} - b^{2}$ is continuous in $t$ and takes the values $1$ and $-1$, so $C_{t}c_{\lambda}$ meets every value of $m$ in the fibre, and an $f$ invariant under the $C_{t}$ is constant on it. A function $V_{0} + k\mu_{Y}$ with $V_{0}$ invariant changes under $C_{2}$ by $-2k\mu_{Y}(c_{\lambda}) = -km(c_{\lambda})$, nonzero for $k \neq 0$ off the balanced frame; the check verifies this, the invariance of $g^{2}\sum\mu_{j}^{2} + g_{R}^{2}\sum\mu_{R,a}^{2}$ (with $g_{R} = \tfrac34$) under four $C_{t}$ at four points, and that the electroweak D-term with level $\zeta = \tfrac52$ changes under $C_{2}$. Clause (2) is the standard fact that a Fayet-Iliopoulos level, a constant shift of a moment map in a D-term, is gauge-invariant only along the centre of the gauge algebra, so it exists only for abelian factors (Fayet and Iliopoulos 1974, Fact \claimref{synthesis.imported-fayet-iliopoulos-level-only-for-abelian-factors}), here made explicit on the rung summand; it is transport-absent. Identification: the four real rung directions span the bidoublet $(2,2,1)$ of $\su(2)_{L} \oplus \su(2)_{R} \oplus \su(4)$ (Proposition \claimref{standard-model.generation-ten-under-pati-salam}), on which the compact $\su(2)_{R}$ preserves $Q$ (it lies in the orthogonal algebra) and $G$ (compact real form) and commutes with $\su(2)_{L}$, which acts as the weak $\su(2)$ of the slot mixing; it is therefore three-dimensional inside the commutant computed above, which it fills. Since $B - L$ vanishes on the bidoublet, the hypercharge acts there as $T_{3R}$, consistent with $Y$ lying in the commutant.

\emph{(3).} $C_{2}$ has $s = 1$, so $C_{2} = 2R_{1}$, and $S(0, \lambda) = (\lambda, 0)$, $S(0, -H^{2}/\lambda) = (-H^{2}/\lambda, 0)$, so $S C_{2}c_{\lambda} = c_{H^{2}/\lambda}$ ($S$ is real, so it acts by the same matrix on $\beta$). This is the involution of Theorem \claimref{synthesis.rung-unit-transverse-to-norm}(3); the check verifies it at $\lambda = \tfrac12, 3, 7$ with $H^{2} = 2$.

Reading (argued, not part of the statement). The symmetry route to the independence of $v/H$ is closed: the boosts $B_{s}$ of Theorem \claimref{synthesis.rung-unit-transverse-to-norm}(1) satisfy $*B_{s} = B_{1/s}*$ for the tower conjugation, and $\mathsf{N}/|Q| = \cosh 2\eta$ on the family distinguishes every frame, so no symmetry of the stored structure moves the frame; the frame is fixed through the level, a free constant. By (1) the level cannot be traded for a constant inside $W(\mathsf{N})$. By (2), at a stage where $\su(2)_{R}$ is a gauge symmetry the level vanishes and an invariant potential on $L$ fixes at most $\mathsf{N}$, leaving the charge alignment undecided (the mirror lifts of (3) are then gauge-equivalent); the electroweak level must then be the $T_{3R}$-odd remnant of the $\su(2)_{R}$ breaking, of size at least $|m^{*}|$ by (1), about $\mathsf{N}^{*} = v^{2}/2$ for $v \gg H$. The up-down split under the conjugate reading needs a $T_{3R}$-odd weighting for the same reason (\claimref{synthesis.up-down-split-under-conjugate-reading-needs-su2r-odd-weighting}). Whether the stored chain's breaking leaves such a remnant, and of what size, is not decided here.
